\section*{Exercises on \S\arabic{exercisegroup}}
\phantomsection
\addcontentsline{toc}{section}{Exercises on \protect\S\arabic{exercisegroup}}

\begin{enumerate}
\def\labelenumi{\arabic{enumi})}
\item
  Let E be a complete normed space over R, F a topological space, J an interval of R not reducing to a single point; let \((t, \xi) \mapsto A(t, \xi)\) be a map from \(J \times F\) into \(\mathcal{L}(E)\) , such that when \(\xi\) tends to \(\xi_{0}\) , \(A(t, \xi)\) tends uniformly to \(A(t, \xi_{0})\) in J. If \(C(t, t_{0}, \xi)\) is the resolvent of the linear equation \(d\mathbf{x}/dt = A(t, \xi)\) . x show that, for every compact interval \(K \subset J\) , \(C(t, t_{0}, \xi)\) tends uniformly to \(C(t, t_{0}, \xi_{0})\) on \(K \times K\) as \(\xi\) tends to \(\xi_{0}\) .
\item
  Let \(t \mapsto A(t)\) be a regulated map from \(J\) into \(\mathcal{L}(E)\) such that, for any two points \(s, t\) of \(J\) , \(A(s)\) and \(A(t)\) commute. Put \(B(t) = \int_{t_0}^{t} A(s) ds\) . Show that the resolvent \(C(t, t_0)\) of the equation \(d\mathbf{x}/dt = A(t)\mathbf{x}\) is equal to \(\exp(B(t))\) .
\end{enumerate}

If \(t \mapsto A_1(t)\) , \(t \mapsto A_2(t)\) are two regulated functions from J into \(\mathcal{L}(\mathrm{E})\) such that, for any two points \(s\) , \(t\) of J, \(A_1(s)\) , \(A_1(t)\) , \(A_2(s)\) , \(A_2(t)\) commute pairwise, show that

\[
\exp \left(\int_ {t _ {0}} ^ {t} \left(A _ {1} (s) + A _ {2} (s)\right) d s\right) = \exp \left(\int_ {t _ {0}} ^ {t} A _ {1} (s) d s\right) \exp \left(\int_ {t _ {0}} ^ {t} A _ {2} (s) d s\right).
\]

\begin{enumerate}
\def\labelenumi{\arabic{enumi})}
\setcounter{enumi}{2}
\tightlist
\item
  Given the two matrices
\end{enumerate}

\[
A = \left( \begin{array}{c c} 0 & 1 \\ 0 & 0 \end{array} \right), \qquad \qquad B = \left( \begin{array}{c c} 0 & 0 \\ 1 & 0 \end{array} \right)
\]

show that \(\exp(A+B)\neq\exp(A)\exp(B)\) .

\begin{enumerate}
\def\labelenumi{\arabic{enumi})}
\setcounter{enumi}{3}
\item
  If \(P\) is any automorphism in \(\mathcal{L}(\mathrm{E})\) show that \(\exp(PAP^{-1}) = P\exp(A)P^{-1}\) for every \(A \in \mathcal{L}(\mathrm{E})\) .
\item
  Show by an example that a linear equation \(d\mathbf{x} / dt = A.\mathbf{x} + e^{\alpha t}\mathbf{p}(t)\) , where \(A \in \mathcal{L}(\mathrm{E})\) is independent of \(t\) and \(\mathbf{p}\) is a polynomial (with coefficients in E) can have an integral which is equal to a polynomial of the same degree as \(\mathbf{p}\) , even when \(\alpha\) is a characteristic root of \(A\) .
\item
  Let \(f(X)\) be a polynomial of degree n with complex coefficients, and let
\end{enumerate}

\[
\frac {1}{f (\mathrm{X})} = \sum_ {j = 1} ^ {q} \sum_ {h = 1} ^ {n _ {j}} \frac {\lambda_ {j h}}{(\mathrm{X} - r _ {j}) ^ {h}}
\]

be the decomposition of the rational fraction \(1 / f(\mathrm{X})\) into simple elements (Alg., VII. 7, n° 2). Show that, for any regulated function \(b\) , the function

\[
t \mapsto \sum_ {j = 1} ^ {q} \sum_ {h = 1} ^ {n _ {j}} \lambda_ {j h} \int_ {t _ {0}} ^ {t} \frac {(t - s) ^ {h - 1}}{(h - 1) !} e ^ {r _ {j} (t - s)} b (s) d s
\]

is an integral of the equation \(f(\mathrm{D})x = b(t)\) .

\begin{enumerate}
\def\labelenumi{\arabic{enumi})}
\setcounter{enumi}{6}
\tightlist
\item
  \begin{enumerate}
  \def\labelenumii{\alph{enumii})}
  \tightlist
  \item
    Let \(t \mapsto A(t)\) be a map from an interval \(J \subset R\) into \(\mathcal{L}(E)\) , such that, for every \(x \in E\) , the map \(t \mapsto A(t).x\) of J into E is continuous. Show that \(t \mapsto \|A(t)\|\) is bounded on every compact interval \(K \subset J\) (use Gen.~Top., IX, p.~194, th. 2). Under these conditions show that, for every point \((t_{0}, \mathbf{x}_{0}) \in \mathbf{J} \times \mathbf{E}\) the equation \(d\mathbf{x}/dt = A(t).\mathbf{x} + \mathbf{b}(t)\) has one and only one solution defined on J and equal to \(x_{0}\) at the point \(t_{0}\) . Further, if one denotes this solution by \(\mathbf{u}(t, t_{0}, \mathbf{x}_{0})\) , the map \(\mathbf{x}_{0} \mapsto \mathbf{u}(t, t_{0}, \mathbf{x}_{0})\) is a bijective bicontinuous linear map \(C(t, t_{0})\) of E onto itself, which satisfies the relations (6) (IV, p.~179) and (7) (IV, p.~181); further, the map \((s, t) \mapsto C(s, t)\) is continuous from \(J \times J\) into \(\mathcal{L}(E)\) .
  \end{enumerate}
\end{enumerate}

\begin{enumerate}
\def\labelenumi{\alph{enumi})}
\setcounter{enumi}{1}
\tightlist
\item
  Take for E the space of sequences \(\mathbf{x}=(x_{n})_{n\in\mathbf{N}}\) of real numbers such that \(\lim_{n\to\infty}x_{n}=0\) , endowed with the norm \(\|x\|=\sup_{n}|x_{n}|\) , under which E is complete. For each \(t\in J=[0,1]\) , let \(A(t)\) be the linear map of E into itself such that
\end{enumerate}

\[
A (t). \mathbf {x} = \left(\frac {1}{1 + n t} x _ {n}\right) _ {n \in \mathbf {N}}.
\]

Show that \(A(t)\) satisfies the conditions of a) but that the map \(t \mapsto A(t)\) of J into \(\mathcal{L}(\mathrm{E})\) is not continuous at the point t = 0, and that the resolvent \(C(t, t_0)\) of the equation \(d\mathbf{x}/dt = A(t)\) . x is such that the map \(t \mapsto C(t, t_0)\) of J into \(\mathcal{L}(\mathrm{E})\) is not differentiable at the point t = 0.

* 8) Let G be a complete normed algebra with unit e over the field R.

\begin{enumerate}
\def\labelenumi{\alph{enumi})}
\item
  Let \(t \mapsto \mathbf{a}(t)\) be a regulated function on an interval \(J \subset R\) with values in G. Show that the integral u of the linear equation \(d\mathbf{x}/dt = \mathbf{a}(t)\mathbf{x}\) which takes the value e at a point \(t_0 \in J\) is invertible on J, and that its inverse is the solution of the equation \(d\mathbf{x}/dt = -\mathbf{x}\mathbf{a}(t)\) (if v is the integral of this last equation which takes the value e at the point \(t_0\) , consider the linear equations satisfied by uv and vu). Deduce that, for every \(x_0 \in G\) , the integral of \(d\mathbf{x}/dt = \mathbf{a}(t)\mathbf{x}\) which takes the value \(x_0\) at the point \(t_0\) is equal to \(\mathbf{u}(t)\mathbf{x}_0\) .
\item
  Let \(\mathbf{a}(t)\) , \(\mathbf{b}(t)\) be two regulated functions on J, and u and v the integrals of the equations \(d\mathbf{x}/dt = \mathbf{a}(t)\mathbf{x}\) , \(d\mathbf{x}/dt = \mathbf{x}\mathbf{b}(t)\) taking the value e at the point \(t_{0}\) . Show that the integral of the equation \(d\mathbf{x}/dt = \mathbf{a}(t)\mathbf{x} + \mathbf{x}\mathbf{b}(t)\) which takes the value \(x_{0}\) at the point \(t_{0}\) is equal to \(\mathbf{u}(t)\mathbf{x}_{0}\mathbf{v}(t)\) .
\item
  Let \(\mathbf{a}(t)\) , \(\mathbf{b}(t)\) , \(\mathbf{c}(t)\) , \(\mathbf{d}(t)\) be four regulated functions on \(\mathbf{J}\) , and \((\mathbf{u}, \mathbf{v})\) a solution of the system of two linear equations
\end{enumerate}

\[
\frac {d \mathbf {x}}{d t} = \mathbf {a} (t) \mathbf {x} + \mathbf {b} (t) \mathbf {y}, \quad \frac {d \mathbf {y}}{d t} = \mathbf {c} (t) \mathbf {x} + \mathbf {d} (t) \mathbf {y}.
\]

Show that, if \(\mathbf{v}\) is invertible on \(J\) , then \(\mathbf{w} = \mathbf{u}\mathbf{v}^{-1}\) is an integral of the equation \(d\mathbf{z}/dt = \mathbf{b}(t) + \mathbf{a}(t)\mathbf{z} - \mathbf{z}\mathbf{d}(t) - \mathbf{z}\mathbf{c}(t)\mathbf{z}\) (the ``Ricatti equation''); and conversely. Deduce that every integral of this last equation on \(J\) , taking the value \(\mathbf{x}_0\) at the point \(t_0\) , is of the form \((A(t)\mathbf{x}_0 + B(t)\mathbf{e})(C(t)\mathbf{x}_0 + D(t)\mathbf{e})^{-1}\) , where \(A(t), B(t), C(t)\) and \(D(t)\) are continuous maps from \(J\) into \(\mathcal{L}(E)\) satisfying the identity \(U.(xy) = (U.x)y\) .

\begin{enumerate}
\def\labelenumi{\alph{enumi})}
\setcounter{enumi}{3}
\tightlist
\item
  Let \(w_{1}\) be an integral of \(d\mathbf{z}/dt = \mathbf{b}(t) + \mathbf{a}(t)\mathbf{z} - \mathbf{z}\mathbf{d}(t) - \mathbf{z}\mathbf{c}(t)\mathbf{z}\) on J; show that if w is another integral of this equation such that \(w - w_{1}\) is invertible on J then \(w - w_{1}\) can be expressed in terms of integrals of the equations \(d\mathbf{x}/dt = -(\mathbf{d} + \mathbf{c}\mathbf{w}_{1})\mathbf{x}\) and \(d\mathbf{x}/dt = \mathbf{x}(\mathbf{a} - \mathbf{w}_{1}\mathbf{c})\) (consider the equation satisfied by \((\mathbf{w} - \mathbf{w}_{1})^{-1}\) , and use b)).
\end{enumerate}

\begin{enumerate}
\def\labelenumi{\arabic{enumi})}
\setcounter{enumi}{8}
\tightlist
\item
  Let \(y_{k}\) ( \(1 \leqslant k \leqslant n\) ) be \(n\) real functions defined on an interval \(I \subset \mathbf{R}\) and having a continuous \((n - 1)^{th}\) derivative.
\end{enumerate}

\begin{enumerate}
\def\labelenumi{\alph{enumi})}
\item
  Show that if the \(n\) functions \(y_{k}\) are linearly dependent then the matrix \((y_{k}^{(h)}(t))_{0\leqslant h\leqslant n - 1, 1\leqslant k\leqslant n}\) has rank \(< n\) at every point \(t\in I\) .
\item
  Conversely, if for every \(t \in I\) the matrix \((y_{k}^{(h)}(t))_{0 \leqslant h \leqslant n-1, 1 \leqslant k \leqslant n}\) has rank \textless{} n, show that in every nonempty open interval \(J \subset I\) there exists a nonempty open interval \(U \subset J\) such that the restrictions of the \(y_{k}\) to U are linearly dependent (if p is the smallest of the numbers q \textless{} n such that the Wronskians of any q of the functions \(y_{k}\) are identically zero on J, consider a point \(a \in J\) where the Wronskian of p - 1 of the functions \(y_{k}\) is not zero, and show that p of the functions \(y_{k}\) are integrals of a linear equation of order p - 1 on a neighbourhood of a).
\item
  Put \(y_{1}(t) = t^{2}\) for \(t \in \mathbf{R}\) , \(y_{2}(t) = t^{2}\) for \(t \geqslant 0\) , \(y_{2}(t) = -t^{2}\) for \(t < 0\) ; show that \(y_{1}\) and \(y_{2}\) have continuous derivatives on \(\mathbf{R}\) and that \(y_{1}y_{2}' - y_{2}y_{1}' = 0\) , but that \(y_{1}\) and \(y_{2}\) are not linearly dependent on \(\mathbf{R}\) .
\end{enumerate}

* 10) Let \(t \mapsto X(t)\) be a map from an interval \(\mathbf{J} \subset \mathbf{R}\) into the space of complex matrices of order \(n\) . Suppose that the derivative \(t \mapsto X'(t)\) exists and is continuous on \(\mathbf{J}\) , and is such that \(X(t)X'(t) = X'(t)X(t)\) for every \(t \in \mathbf{J}\) .

\begin{enumerate}
\def\labelenumi{\alph{enumi})}
\item
  Suppose further that for each \(t \in J\) the eigenvalues of \(X(t)\) are distinct. Show that then there exists a constant invertible matrix \(P_{0}\) such that \(P_{0}X(t)P_{0}^{-1} = D(t)\) , where \(D(t)\) is a diagonal matrix, so that \(X(t_{1})\) and \(X(t_{2})\) commute for every pair of values \(t_{1}, t_{2}\) of t in J. (Write \(X(t) = P(t)D(t)P(t)^{-1}\) on a neighbourhood of each point of J and form the differential equation satisfied by \(P(t)\) .)
\item
  The matrix
\end{enumerate}

\[
X (t) = \left( \begin{array}{c c c} t ^ {2} & t ^ {3} & t ^ {4} \\ - 2 t & - 2 t ^ {2} & - 2 t ^ {3} \\ 1 & t & t ^ {2} \end{array} \right)
\]

commutes with its derivative, but the conclusion of a) does not hold.*

